% Generated by analysis/build-paper.mjs. Item-level accuracy with 95% Wilson intervals.
\begin{table*}[t]
\caption{Item-level accuracy, 00: A simple command.}\label{tab:items-normal}
\small
\begin{tabular}{p{7.2cm}rr}
\toprule
Item & Correct & 95\% interval\\
\midrule
Which of these happen between pressing Send and the rover moving? Select all that app... (core) & -- & --\\
The natural-language sentence is understood on the rover's board (the ESP32). & -- & --\\
A control program is produced for the command each time it is sent. & -- & --\\
The rover is told what it can do through a registered list of its hardware capabilities. & -- & --\\
\bottomrule
\end{tabular}
\end{table*}

\begin{table*}[t]
\caption{Item-level accuracy, WM2: The rover did not move.}\label{tab:items-wm2}
\small
\begin{tabular}{p{7.2cm}rr}
\toprule
Item & Correct & 95\% interval\\
\midrule
At which point in the process do you think the command stopped? (core) & -- & --\\
Which explanation best fits why the rover did not move? (core) & -- & --\\
If the rover does not move, the motors must be faulty. & -- & --\\
In this attempt the rover's board was flashed with a new program. & -- & --\\
Different kinds of fault can look identical from the outside: the rover just doesn't ... & -- & --\\
Replacing the motors would fix this failure. & -- & --\\
The same command can succeed later once the underlying fault is removed. & -- & --\\
\bottomrule
\end{tabular}
\end{table*}

\begin{table*}[t]
\caption{Item-level accuracy, WM3: Find the red ball.}\label{tab:items-wm3}
\small
\begin{tabular}{p{7.2cm}rr}
\toprule
Item & Correct & 95\% interval\\
\midrule
Could this rover find a red ball if the command were phrased differently, or the AI t... (core) & -- & --\\
What the rover can do is limited by the hardware registered on its board, not only by... & -- & --\\
A more capable AI model would let the rover see without a camera. & -- & --\\
The system generated motion code for this request, tried it, and failed to find the b... & -- & --\\
The refusal means the rover is broken. & -- & --\\
This rover has no camera, microphone, distance sensor or gripper. & -- & --\\
\bottomrule
\end{tabular}
\end{table*}

\begin{table*}[t]
\caption{Item-level accuracy, WM4: ``Do that again''.}\label{tab:items-wm4}
\small
\begin{tabular}{p{7.2cm}rr}
\toprule
Item & Correct & 95\% interval\\
\midrule
Why did ``Do that again'' fail? (core) & -- & --\\
Where was the record of your first instruction held? (core) & -- & --\\
Restating the whole command (``Move forward for 2 seconds'') would work. & -- & --\\
The follow-up failed because the motors had a fault. & -- & --\\
``Do that again'' only works while an earlier action is still in the agent's active con... & -- & --\\
After this failure the rover has permanently lost the ability to repeat actions. & -- & --\\
The rover chose to ignore the request. & -- & --\\
\bottomrule
\end{tabular}
\end{table*}

\begin{table*}[t]
\caption{Item-level accuracy, WM6: Who does the thinking?.}\label{tab:items-wm6}
\small
\begin{tabular}{p{7.2cm}rr}
\toprule
Item & Correct & 95\% interval\\
\midrule
Where is ``Turn left'' turned into something the rover can execute? (core) & -- & --\\
The rover was connected, yet nothing happened. Which explanation fits best? (core) & -- & --\\
A connected board can interpret any new English request on its own. & -- & --\\
Restoring the agent connection, not repairing the rover, is what was needed. & -- & --\\
No new program was generated in this attempt. & -- & --\\
The rover ``decided'' not to turn. & -- & --\\
\bottomrule
\end{tabular}
\end{table*}

\begin{table*}[t]
\caption{Item-level accuracy, WM7: It used to work.}\label{tab:items-wm7}
\small
\begin{tabular}{p{7.2cm}rr}
\toprule
Item & Correct & 95\% interval\\
\midrule
The same command worked earlier and now fails. Which explanation fits best? (core) & -- & --\\
A rule in the system can change even though I did nothing differently. & -- & --\\
Sending the identical 4-second command again will eventually succeed. & -- & --\\
A shorter command (e.g. 2 seconds) would still be allowed. & -- & --\\
The 4-second program was flashed and then stopped part-way. & -- & --\\
\bottomrule
\end{tabular}
\end{table*}

\begin{table*}[t]
\caption{Item-level accuracy, LEARN: After the fix.}\label{tab:items-ls}
\small
\begin{tabular}{p{7.2cm}rr}
\toprule
Item & Correct & 95\% interval\\
\midrule
After the fix, the same command worked. Why? (core) & -- & --\\
The rover now remembers this failure and will change its future behaviour. & -- & --\\
A new session starts without the memory of this session. & -- & --\\
The rover's skills improve the more I use it. & -- & --\\
Repeated successful use can change what hardware capabilities the rover has. & -- & --\\
\bottomrule
\end{tabular}
\end{table*}

\begin{table*}[t]
\caption{Item-level accuracy, WM5: Vignette: early success.}\label{tab:items-wm5}
\small
\begin{tabular}{p{7.2cm}rr}
\toprule
Item & Correct & 95\% interval\\
\midrule
Is Priya's expectation well-founded? (core) & -- & --\\
After five successes, the chance of any failure is effectively zero. & -- & --\\
A failure after several successes means the robot has become unreliable overall. & -- & --\\
A task can fail because it needs a capability the robot never had, however many other... & -- & --\\
\bottomrule
\end{tabular}
\end{table*}

\begin{table*}[t]
\caption{Item-level accuracy, WM1: Vignette: two languages.}\label{tab:items-wm1}
\small
\begin{tabular}{p{7.2cm}rr}
\toprule
Item & Correct & 95\% interval\\
\midrule
Which explanation fits best? (core) & -- & --\\
The failure necessarily means the motors changed. & -- & --\\
The same hardware can work well in one language and poorly in another. & -- & --\\
\bottomrule
\end{tabular}
\end{table*}

\begin{table*}[t]
\caption{Item-level accuracy, OVERALL: Putting it together.}\label{tab:items-arch}
\small
\begin{tabular}{p{7.2cm}rr}
\toprule
Item & Correct & 95\% interval\\
\midrule
Interprets the meaning of the English sentence (core) & -- & --\\
Writes the control program (core) & -- & --\\
Keeps track of earlier commands (context) (core) & -- & --\\
Checks the program before it is sent to the rover (core) & -- & --\\
Physically drives the motors when the program runs (core) & -- & --\\
The rover thinks and decides on its own what to do. & -- & --\\
There are several separate places a command can fail before the rover moves. & -- & --\\
A failure to move always points to a hardware problem. & -- & --\\
What the rover can do is determined by installed hardware and firmware limits, not by... & -- & --\\
\bottomrule
\end{tabular}
\end{table*}
